\documentclass[a4paper, 11pt, oneside, oldfontcommands]{memoir}

%%%%% Packages %%%%%
\usepackage{lmodern}
\usepackage{palatino}
\usepackage[T1]{fontenc}
\usepackage[utf8]{inputenc}
\usepackage[english]{babel}


%%%%%%%%%%%%%%%%%%%%  PACKAGE SECONDAIRE

%\usepackage{amstext,amsmath,amssymb,amsfonts} % package math
%\usepackage{multirow,colortbl}	% to use multirow and ?
%\usepackage{xspace,varioref}
\usepackage[linktoc=all, hidelinks]{hyperref}			% permet d'utiliser les liens hyper textes
\usepackage{float}				% permet d ajouter d autre fonction au floatant
%\usepackage{wrapfig}			% permet d avoir des image avec texte coulant a cote
%\usepackage{fancyhdr}			% permet d inserer des choses en haut et en bas de chaque page
\usepackage{microtype}			% permet d ameliorer l apparence du texte
%\usepackage[explicit]{titlesec}	% permet de modifier les titres
\usepackage{graphicx}			% permet d utiliser les graphiques
\graphicspath{{./images/}}		% to say where are image
%\usepackage{eso-pic} 			% to put figure in the background
\usepackage[svgnames]{xcolor}	% permet d avoir plus de 300 couleur predefini
%\usepackage{array}				% permet d ajouter des option dans les tableaux
%\usepackage{listings}			% permet d ajouter des ligne de code
%\usepackage{tikz}				% to draw figure
%\usepackage{appendix}			% permet de faire les index
%\usepackage{makeidx}			% permet de creer les index
%\usepackage{fancyvrb}			% to use Verbatim
%\usepackage{framed}				% permet de faire des environnement cadre
%\usepackage{fancybox}			% permet de realiser les cadres
\usepackage{titletoc}			% permet de modifier les titres
%\usepackage{caption}
%\usepackage[a4paper, top=2cm, bottom=2cm]{geometry}
\usepackage{frbib}                      %permet d avoir une biblio francaise
\usepackage[babel=true]{csquotes}


\usepackage{graphicx}
\RequirePackage{pageGardeEnsta}	% permet d avoir la page de garde ensta

\setcounter{secnumdepth}{2}		% permet d'augmenter la numerotation
%\setcounter{tocdepth}{2}		% permet d'augmenter la numerotation

%%%%%%%%%%%%%%%%%%  DEFINITION DES BOITES
\newcounter{rem}[chapter]

\newcommand{\remarque}[1]{\stepcounter{rem}\noindent\fcolorbox{OliveDrab}{white}{\parbox{\textwidth}{\textcolor{OliveDrab}{
\textbf{Remarque~\thechapter.\therem~:}}\\#1}}}

\newcounter{th}[chapter]

\newcommand{\theoreme}[2]{\noindent\fcolorbox{FireBrick}{white}{\stepcounter{th}
\parbox{\textwidth}{\textbf{\textcolor{FireBrick}{Théorème~\thechapter.\theth~:}}{\hfill \textit{#1}}\\#2}}}

\newcommand{\attention}[1]{\noindent\fcolorbox{white}{white}{\parbox{\textwidth}{\textcolor{FireBrick}{
        \textbf{Attention !}}\\\textit{#1}\\}}}


\newcommand{\definition}[2]{\noindent\fcolorbox{OliveDrab}{white}{\stepcounter{th}
    \parbox{\textwidth}{\textbf{\textcolor{OliveDrab}{Définition~\thechapter.\theth~:}}{~
        \textit{#1}}~\setlength{\parindent}{15pt}\par#2\par}}\par~\\}


%%%%%%%%%%%%%%%%%%%%%%%%%%%%%%%%%%%%%%%%%%%%%%%%%%%%%%%%%%%%%%%%%%%%%%%%%


%% INDEX %%%%%%%%%%%%%%%%%%%%%%%%%%%%%%%%%%%%%%%%%%%%%%%%%%%%
\makeindex

%%%%% Useful macros %%%%%
\newcommand{\latinloc}[1]{\ifx\undefined\lncs\relax\emph{#1}\else\textrm{#1}\fi\xspace}
\newcommand{\etc}{\latinloc{etc}}
\newcommand{\eg}{\latinloc{e.g.}}
\newcommand{\ie}{\latinloc{i.e.}}
\newcommand{\cad}{c'est-à-dire }
\newcommand{\st}{\ensuremath{\text{\xspace s.t.\xspace}}}

%%%% Definition des couleur %%%%

\newcommand\couleurb[1]{\textcolor{SteelBlue}{#1}}
\newcommand\couleurr[1]{\textcolor{DarkRed}{#1}}


%% number page style style %%%%%%%%%%%%%%%%%%%%%%%%%%%%%%%%%%%%%%%%%%%%%%%%%%%%%%

\pagestyle{plain}
%\pagestyle{empty}
%\pagestyle{headings}
%\pagestyle{myheadings}



%% chapters style %%%%%%%%%%%%%%%%%%%%%%%%%%%%%%%%%%%%%%%%%%%%%%%%%%%%%%
%% You may try several styles (see more in the memoir manual).

%\chapterstyle{veelo}
%\chapterstyle{chappell}
%\chapterstyle{ell}
%\chapterstyle{ger}
%\chapterstyle{pedersen}
%\chapterstyle{verville}
\chapterstyle{madsen}
%\chapterstyle{thatcher}


%%%%% Report Title %%%%%
\title{Sniff Hynesim}
\author{\textsc{Rigaud Michaël}}
\date{\today}
\doctype{Status Report}
\promo{promo 2017}
\etablissement{\textsc{Ensta} Bretagne\\2, rue François Verny\\
  29806 \textsc{Brest} cedex\\\textsc{France}\\Tel +33 (0)2 98 34 88 00\\ \url{www.ensta-bretagne.fr}}
\logoEcole{\includegraphics[height=4.2cm]{logo_ENSTA_Bretagne_Vertical_CMJN}}



%%%%%%%%%%%%%%%%%% DEBUT DU DOCUMENT
\begin{document}

\maketitle
\thispagestyle{empty}
\newpage

\tableofcontents


%%%%%%%%%%%%%%%%% INTRODUCTION

\chapter*{Introduction}
\addcontentsline{toc}{chapter}{Introduction}

The cyber security is one of the major thread of the 21th century, and attackers
use techniques more and more sophisticated. So one of the most important aim for
cyber security engineer is to find a way to detect and stop attacks. In this
project we decide to elaborate a solution to alert when a strategy of attack is
spotted out. To do that, we have to create pattern of attack and use an
IDS\footnote{Intrusion detection system} to alert us. To verify the solution and
create pattern as exhaustive as possible, we decide to use Hynesim. It is a
solution of network virtualization which permit a huge agility.

To begin, we will present Hynesim and the advantages of tis software. Then, we
are going to present the aim of an IDS and the most popular IDS. And to finish
we will present the aim of this project.
\newpage
%%%%%%%%%%%%%%%%%%%%%%%%


\chapter{Subject}

\section{Context}

\section{Aim}



%%% Local Variables:
%%% mode: latex
%%% TeX-master: "../rapport_de_base"
%%% End:


\chapter{Hynesim}
\label{chap:hynesim}

\section{Presentation}
\definition{Hynesim}{Means HYbrid NEtwork SIMulation, is a distribute platform
  of simulation of information system developed by Diateam. \cite{hynesim}}

The platform was initially developed by Diateam for DGA MI (Maitrise de
l'information) to create virtual network. But now is a major project to develop
information system and automatize cyber security attacks. This project has two
version, an open source version and an professional version. The open source
version as less option, but we will use this version for this project.

\section{Architecture}

To work, Hynesim need a server with on it the main software. This software is
the virtualization part. It manage virtual machine and network.

Moreover, to see virtual machine and interact with them, users need to have the
client interface. This interface can be install on a simple computer.

%%%%%%%%%%%%%%%%%%%%%%%%%%%%%%%%%%%%%%%%%%%%%%%%%%%%%%%%%%%%%%%%%%%%%%%%%%%%%%%
%% TODO: Ajouter image d un reseau hynesim
%%%%%%%%%%%%%%%%%%%%%%%%%%%%%%%%%%%%%%%%%%%%%%%%%%%%%%%%%%%%%%%%%%%%%%%%%%%%%%%




%%% Local Variables:
%%% mode: latex
%%% TeX-master: "../rapport_de_base"
%%% End:

\chapter{IDS}

In this subject, we will have to create probe to detect attack. It is the aim of IDS that we will present.

\section{Presentation}

\definition{IDS}{An intrusion detection system (IDS) inspects all inbound and outbound network activity and
identifies suspicious patterns that may indicate a network or system attack from someone attempting to break into
or compromise a system. \cite{webopedia}}



There is many type of IDS:
\begin{itemize}[itemsep=0pt]
\item NIDS, network IDS. They listen and analyze the network and detect attack from network packet. They are the
  most interesting for our subject, so this document will focus primarily on this type of IDS.
\item HIDS, host IDS. These IDS are on a system and they detect intrusion inside it.
\item Hybrid IDS. They are composed with NIDS and HIDS.
\item IPS. They are NIDS with active functions which permit to stop attackers.
\item KIDS/KIPS, kernel IDS. They are type of HIDS. They are on the system kernel. They are more effective and
slower than HIDS.
\end{itemize}

In the following of this document we will talk about NIDS. But to detect attack there is also many methods.
~\\


To be efficient and IDS should have a good balance between some features.
\begin{itemize}[itemsep=0pt]
\item \textbf{Speed:} In fact an IDS should analyze packet as fast as possible, otherwise, it will behind the network
  traffic.
\item \textbf{False alarm:} An IDS raise an alert when it detect attack. But it could raise alarm during a normal
  utilization. It is name false alarm. It is one of the most important feature, because at every alarm a system
  administrator need to analyze the alert. So in company, every alarm cost time and money.
\item \textbf{Probability of detection:} It is for an IDS the capacity of detecting attack. More this probability is high
  more the IDS will raise false alarm. In some sensitive system, we prefer detect every attack and raise many
  false alarm. That depends on the system.
\end{itemize}


\section{Detection methods}

\subsection{Misuse detection}

This technique is the simplest. It use attack signature to raise alert. In fact, all attacks have a particularity,
if we detect this particularity we can detect the attackers. There is three sub methods.

\subsubsection{Pattern matching}

In this technique we have a based of signature and the IDS is looking for the pattern. If the pattern match
perfectly, this IDS raise an alert.

The problem is, only attacks which are in our based can be detected. So if there is a new attack (zero day), or if
the attack is not perfectly the same, we can't detect it.

However, this method is much used because it is high-performance, and with this method the IDS don't raise a lot of
false alerts

\subsubsection{Dynamic pattern matching}

In this techniques the IDS is also based on signature but this data base is dynamic. In other words, the IDS has
the faculty of adaptation and learning. The IDS improve his data base of signature automatically.

\subsubsection{Protocol analysis}

The last sub method we will present is the protocol analysis. This technique is based on the verification of
protocol. The IDS will check if flows are compliant with RFC\footnote{Requests for Comments, is a type of
  publication from the Internet Engineering Task Force (IETF) and the Internet Society (ISOC), the principal
  technical development and standards-setting bodies for the Internet.} standards. It will verify parameters of
packets and fields of them. An IDS can check many protocol as FTP, HTTP, ICMP, \dots

The advantages of this methods is that we can detect unknown attacks in contrary of pattern matching. However,
software publishers don't often respect RFC so this technique is not always very efficient.





\subsection{Anomaly detection}

This technique consists in detecting an intrusion with the analysis of the user's past behavior. So the IDS should
create a profile of users from his use and raise alert when there is an event outside this profile. To create
profile the IDS could use machine learning.
~\\

\begin{tabular}{|p{0.45\textwidth}|p{0.45\textwidth}|} \hline
Advantages                                                                                                 & Drawbacks \\ \hline The IDS should be able to detect every type of
attack even unknown (zero days) attacks.                                                                   & This method is not reliable. Every
alteration of the use create an alert                                                                                  \\
  \hline The IDS is autonomous                                                                             & This method
need a learning period. In fact, this method need to learn the habit of users,
so we need a period without attacks                                                                                    \\ \hline & An hackers can need only time. In fact,
if the attacker arrive to create a new profile after month with his attacks, he
could attack silently                                                                                                  \\ \hline

\end{tabular}


\subsubsection{Probabilistic method}
Bayésien network is a learning machine based on probability. The IDS will create a probabilistic model and if the
user and will raise an alert if the user don't respect this model.

For example, we know that in 90\% of cases in an HTTP request the first parameter is GET after a connection to the
port 80.


%
% \subsubsection{Statistic method}
% \subsubsection{Other method}

% It is possible to use every type of machine learning. In particularly, neural
% network.


\section{Main IDS}

There is many IDS on the market. The most popular open source solution are Snort, Suricata, Bro, Fail2Ban, ACARM
... There is also some tools <<all in one>>. It is usually OS\footnote{Operating system} with an IDS, a tool to
analyze alert, a tool to create rules,... The most popular are SELKS and ELKS. A description of SELKS is available
page \pageref{chap:selks}.




%%% Local Variables:
%%% mode: latex
%%% TeX-master: "../rapport_de_base"
%%% End:



%%%% CONCLUSION %%%%%%%%%

\chapter*{Conclusion}
\addcontentsline{toc}{chapter}{Conclusion}
\newpage

%%%% ANNEXE %%%%%%%%%%%%

%\part*{Annexe}
\appendix
\nocite{*}
%\input{annexe_}
\newpage
 \listoffigures
 \printindex
 \bibliographystyle{plain}
  \bibliography{biblio}

\end{document}
%%%%%%%%%%%%%%%%% FIN DU DOCUMENT
%%% Local Variables:
%%% mode: latex
%%% TeX-master: t
%%% End:
